\begin{fortranspecific}[4ex]
\section{\kcode{workdistribute} Construct}
\label{sec:target_teams_workdistribute}
\index{constructs!workdistribute@\kcode{workdistribute}}
\index{workdistribute construct@\kcode{workdistribute} construct}

The following examples illustrate the use of the combined \kcode{target teams
workdistribute} construct to execute and parallelize Fortran array statements on
a device.  All examples may alternatively be expressed as a \kcode{workdistribute}
construct closely nested inside a \kcode{target teams} construct.


In the following example, \kcode{target teams workdistribute} is used to
evaluate the \emph{axpy} formula. The individual element-wise operations are
automatically assigned to the OpenMP threads on the device such that the
array operations run in parallel across the league of teams.

\ffreenexample{target_teams_workdistribute}{1}

\topmarker{Fortran}
The next example shows a \kcode{target teams workdistribute} construct that
contains several array statements. The statements have dependences between
them and the implementation will ensure that the array statements are exeuted
in the correct sequence. Therefore, the implementation may introduce additional
synchronization between the statements or may split the construct into several
\kcode{target teams} constructs to maintain Fortran semantics.

\ffreenexample{target_teams_workdistribute}{2}

The last example uses \kcode{target teams workdistribute} to perform a more
complex computation. The computation of \ucode{minval(dd)} may be
parallelized while the assignment to scalar variable \ucode{f} is performed only by a
single thread.

\ffreenexample{target_teams_workdistribute}{3}

\end{fortranspecific}
